\honest{Applause Lights}{Eliezer Yudkowsky}{2007}

At the 2007 Singularity Summit, a speaker called for democratic, multinational development of artificial intelligence. I went to the microphone and asked which was more ``democratic'': democracies forming a consortium full of ordinary politicking, or rebel nerds whose AI polls the whole world and does what the majority says. I wanted to learn whether he valued democracy for how well it works or for the moral rightness of voting. The speaker said the first sounded like an editorial in \textsc{Reason} and the second like a Hollywood plot. \nb{The speaker, Jamais Cascio, wrote in the comments that the reply went on: both scenarios were ``caricatures,'' each with benefits and drawbacks.} I asked what kind of democratic process he had in mind, and he said something like the Human Genome Project. \nb{Pointing readers to the recording, Yudkowsky wrote in the comments that Cascio had offered the Genome Project example without being asked, and that the event ``did not occur the exact way I remembered it.'' The post as published today still shows the example as an answer to the question.} I asked how interest groups would resolve their conflicts in such a structure. He said, ``I don't know.'' \nb{By Cascio's account the answer continued: ``this is not a solved problem, but it's an important problem, and we need to figure out how to address it.''}

This reminded me of ``some dictator or other'' who said his state was already democratic, although ``the expression of the people's will'' was still missing. \nb{It was Roberto Eduardo Viola of Argentina, in a \textsc{Time} interview in 1981.}

The substance of a democracy is a mechanism for resolving policy conflicts. If everyone wanted the same policies, no democracy would be needed. The mechanism can be a majority vote, an elected legislature or even an AI that responds to voters, but it has to be something. Without one in mind, a call for democracy means only that ``you have said the word `democracy,' so the audience is supposed to cheer.'' It is an applause light. \nb{George Orwell made this point about the word ``democracy'' in ``Politics and the English Language'' in 1946. Yudkowsky said in the comments that Orwell's essay was in mind while writing; the post does not mention it.}

Most applause lights are easier to spot than this one. Reverse the sentence. ``We shouldn't balance the risks and opportunities of AI'' sounds abnormal, so the original ``probably'' tells you nothing new. Such sentences can introduce a topic or stress a proposal, but ``if no specifics follow, the sentence is probably an applause light.'' \nb{A specific did follow in the Summit exchange: the Human Genome Project. Yudkowsky allowed in the comments, ``I do not accuse you too much of this.''}

I end with a talk I am tempted to give, made of nothing but applause lights, to see how long the audience takes to start laughing: we must balance the risks and opportunities, plan wisely and rationally, avoid both fear and blind enthusiasm, respect all stakeholders, share the benefits widely, and think before it is too late.
